\SourceLeaf{087}
polygona formantur, in triangula per mensurationes non divisa. Sed de his rebus, ab instituto praesenti nimis alienis, alia occasione fusius agemus. Silentio tamen praeterire non possumus monitum, quod theoria nostra, si applicatio pura atque rigorosa in votis est, supponit, quantitates per \(v,v',v''\) etc. designatas revera vel immediate observatas esse, vel ex observationibus ita derivatas, ut inter se independentes maneant, vel saltem tales censeri possint. In praxi vulgari observantur anguli triangulorum ipsi, qui proin pro \(v,v',v''\) etc. accipi possunt; sed memores esse debemus, si forte systema insuper contineat triangula talia, quorum anguli non sint immediate observati, sed prodeant tamquam summae vel differentiae angulorum revera observatorum, illos non inter observatorum numerum referendos, sed in forma compositionis suae in calculis retinendos esse. Aliter vero res se habebit in modo observandi ei simili, quem sequutus est clar. \textsc{Struve} (Astronomische Nachrichten II, p. 431), ubi directiones singulorum laterum ab eodem vertice proficiscentium obtinentur per comparationem cum una eademque directione arbitraria. Tunc scilicet hi ipsi anguli pro \(v,v',v''\) etc. accipiendi sunt, quo pacto omnes anguli triangulorum in forma differentiarum se offerent, aequationesque conditionales primi generis, quibus per rei naturam sponte satisfit, tamquam superfluae cessabunt. Modus observationis, quem ipse sequutus sum in dimensione triangulorum annis praecedentibus perfecta, differt quidem tum a priore tum a posteriore modo, attamen respectu effectus posteriori aequiparari potest, ita ut in singulis stationibus directiones laterum inde proficiscentium ab initio quasi arbitrario numeratas pro quantitatibus \(v,v',v''\) etc. accipere oporteat. Duo iam exempla elaborabimus, alterum ad modum priorem, alterum ad posteriorem pertinens.

\Article{23}
Exemplum primum nobis suppeditabit opus clar. \textsc{de Krayenhof}, \textit{Précis historique des opérations trigonométriques faites en Hollande}, et quidem compensationi subiiciemus partem eam systematis triangulorum, quae inter novem puncta Harlingen, Sneek, Oldeholtpade, Ballum, Leeuwarden, Dockum, Drachten, Oosterwolde, Gröningen continetur. Formantur inter haec puncta novem triangula in opere illo per numeros 121, 122, 123, 124, 125, 127, 128, 131, 132 denotata, quorum anguli (a nobis indicibus praescriptis distincti) secundum tabulam p. 77---81 ita sunt observati:
\SourceLeaf{088}
\begin{center}
\begin{tabular}{r l l r r r}
\multicolumn{6}{c}{Triangulum 121.}\\[5pt]
0.&Harlingen&. . . . .&\(50^0\)&\(58'\)&\(15''238\)\\
1.&Leeuwarden&. . .&82&47&15,351\\
2.&Ballum&. . . . . . .&46&14&27,202
\end{tabular}
\end{center}
\begin{center}
\begin{tabular}{r l l r r r}
\multicolumn{6}{c}{Triangulum 122.}\\[5pt]
3.&Harlingen&. . . . .&51&5&39,717\\
4.&Sneek&. . . . . . . .&70&48&33,445\\
5.&Leeuwarden&. . .&58&5&48,707
\end{tabular}
\end{center}
\begin{center}
\begin{tabular}{r l l r r r}
\multicolumn{6}{c}{Triangulum 123.}\\[5pt]
6.&Sneek&. . . . . . . .&49&30&40,051\\
7.&Drachten&. . . . . .&42&52&59,382\\
8.&Leeuwarden&. . .&87&36&21,057
\end{tabular}
\end{center}
\begin{center}
\begin{tabular}{r l l r r r}
\multicolumn{6}{c}{Triangulum 124.}\\[5pt]
9.&Sneek&. . . . . . . .&45&36&7,492\\
10.&Oldeholtpade&. . .&67&52&0,048\\
11.&Drachten&. . . . .&66&31&56,513
\end{tabular}
\end{center}
\begin{center}
\begin{tabular}{r l l r r r}
\multicolumn{6}{c}{Triangulum 125.}\\[5pt]
12.&Drachten&. . . . . .&53&55&24,745\\
13.&Oldeholtpade&. . .&47&48&52,580\\
14.&Oosterwolde&. . .&78&15&42,347
\end{tabular}
\end{center}
\begin{center}
\begin{tabular}{r l l r r r}
\multicolumn{6}{c}{Triangulum 127.}\\[5pt]
15.&Leeuwarden&. . .&59&24&0,645\\
16.&Dockum&. . . . . .&76&34&9,021\\
17.&Ballum&. . . . . . .&44&1&51,040
\end{tabular}
\end{center}
\begin{center}
\begin{tabular}{r l l r r r}
\multicolumn{6}{c}{Triangulum 128.}\\[5pt]
18.&Leeuwarden&. . .&72&6&32,043\\
19.&Drachten&. . . . . .&46&53&27,163\\
20.&Dockum&. . . . . .&61&0&4,494
\end{tabular}
\end{center}
\begin{center}
\begin{tabular}{r l l r r r}
\multicolumn{6}{c}{Triangulum 131.}\\[5pt]
21.&Dockum&. . . . . .&57&1&55,292\\
22.&Drachten&. . . . . .&83&33&14,515\\
23.&Gröningen&. . . . .&39&24&52,397
\end{tabular}
\end{center}
\SourceLeaf{089}
\begin{center}
\begin{tabular}{r l l r r r}
\multicolumn{6}{c}{Triangulum 132.}\\[5pt]
24.&Oosterwolde&. . .&\(81^0\)&\(54'\)&\(17''447\)\\
25.&Gröningen&. . . . .&31&52&46,094\\
26.&Drachten&. . . . .&66&12&57.246
\end{tabular}
\end{center}

Consideratio nexus inter haec triangula monstrat, inter 27 angulos, quorum valores approximati per observationem innotuerunt, 13 aequationes conditionales haberi, puta duas primi generis, novem secundi, duas tertii. Sed haud opus erit, has aequationes omnes in forma sua finita hic adscribere, quum ad calculos tantummodo requirantur quantitates in theoria generali per \(\mathfrak A,a,a',a''\) etc., \(\mathfrak B,b,b',b''\) etc. etc. denotatae: quare illarum loco statim adscribimus aequationes supra per (13) denotatas, quae illas quantitates ob oculos ponunt: loco signorum \(\varepsilon,\varepsilon',\varepsilon''\) etc. simpliciter hic scribemus \((0),(1).\ (2)\) etc.

Hoc modo duabus aequationibus conditionalibus primi generis respondent sequentes:
\[
\begin{aligned}
(1)+(5)+(8)+(15)+(18)&=-2''197\\
(7)+(11)+(12)+(19)+(22)+(26)&=-0''436
\end{aligned}
\]

Excessus sphaeroidicos novem triangulorum invenimus deinceps: \(1''749;\allowbreak1''147;\allowbreak1''243;\allowbreak1''698;\allowbreak0''873;\allowbreak1''167;\allowbreak1''104;\allowbreak2''161;\allowbreak1''403\). Oritur itaque aequatio conditionalis secundi generis prima haec\footnote{. Indices in hoc exemplo per figuras arabicas exprimere praeferimus.}: \(v^{(0)}+v^{(1)}+v^{(2)}-180^0\,0'\,1''749=0\), et perinde reliquae: hinc habemus novem aequationes sequentes:
\[
\begin{aligned}
(0)+(1)+(2)&=-3''958\\
(3)+(4)+(5)&=+0{,}722\\
(6)+(7)+(8)&=-0{,}753\\
(9)+(10)+(11)&=+2{,}355\\
(12)+(13)+(14)&=-1{,}201\\
(15)+(16)+(17)&=-0{,}461\\
(18)+(19)+(20)&=+2{,}596\\
(21)+(22)+(23)&=+0{,}043\\
(24)+(25)+(26)&=-0{,}616
\end{aligned}
\]
Aequationes conditionales tertii generis commodius in forma logarithmica exhibentur: ita prior est
